\section{Localization}
\label{section-localize}

\noindent
Let $\mathcal{C}$ be a site.
Let $U \in \Ob(\mathcal{C})$.
See
Categories, Example \ref{categories-example-category-over-X}
for the definition of the category $\mathcal{C}/U$ of
objects over $U$. We turn $\mathcal{C}/U$ into a site
by declaring a family of morphisms $\{V_j \to V\}$ of
objects over $U$ to be a covering of $\mathcal{C}/U$
if and only if it is a covering in $\mathcal{C}$.
Consider the forgetful functor
$$
j_U : \mathcal{C}/U \longrightarrow \mathcal{C}.
$$
This is clearly cocontinuous and continuous. Hence by
Lemma \ref{lemma-when-shriek} we obtain a morphism of topoi
$$
j_U : \Sh(\mathcal{C}/U) \longrightarrow \Sh(\mathcal{C})
$$
given by $j_U^{-1}$ and $j_{U*}$ and a functor $j_{U!}$
left adjoint to $j_U^{-1}$.

\begin{definition}
\label{definition-localize}
Let $\mathcal{C}$ be a site.
Let $U \in \Ob(\mathcal{C})$.
\begin{enumerate}
\item The site $\mathcal{C}/U$ is called the {\it localization of
the site $\mathcal{C}$ at the object $U$}.
\item The morphism of topoi
$j_U : \Sh(\mathcal{C}/U) \to \Sh(\mathcal{C})$
is called the {\it localization morphism}.
\item The functor $j_{U*}$ is called the {\it direct image functor}.
\item For a sheaf $\mathcal{F}$ on $\mathcal{C}$ the sheaf
$j_U^{-1}\mathcal{F}$ is called the {\it restriction of $\mathcal{F}$
to $\mathcal{C}/U$}.
\item For a sheaf $\mathcal{G}$ on $\mathcal{C}/U$
the sheaf $j_{U!}\mathcal{G}$ is called the
{\it extension of $\mathcal{G}$ by the empty set}.
\end{enumerate}
\end{definition}

\noindent
The restriction $j_U^{-1}\mathcal{F}$ is the sheaf
defined by the rule $j_U^{-1}\mathcal{F}(X/U) = \mathcal{F}(X)$ as expected.
The extension by the empty set also has a very easy description in
this case; here it is.

\begin{lemma}
\label{lemma-describe-j-shriek}
Let $\mathcal{C}$ be a site.
Let $U \in \Ob(\mathcal{C})$.
Let $\mathcal{G}$ be a presheaf on $\mathcal{C}/U$.
Then $j_{U!}(\mathcal{G}^\#)$ is the sheaf associated to the presheaf
$$
V
\longmapsto
\coprod\nolimits_{\varphi \in \Mor_\mathcal{C}(V, U)}
\mathcal{G}(V \xrightarrow{\varphi} U)
$$
with obvious restriction mappings.
\end{lemma}

\begin{proof}
By Lemma \ref{lemma-when-shriek} we have
$j_{U!}(\mathcal{G}^\#) = ((j_U)_p\mathcal{G}^\#)^\#$.
By Lemma \ref{lemma-technical-up} this is equal to $((j_U)_p\mathcal{G})^\#$.
Hence it suffices to prove that $(j_U)_p$ is given by
the formula above for any presheaf $\mathcal{G}$ on $\mathcal{C}/U$.
OK, and by the definition in Section \ref{section-functoriality-PSh} we have
$$
(j_U)_p\mathcal{G}(V)
=
\colim_{(W/U, V \to W)} \mathcal{G}(W)
$$
Now it is clear that the category of pairs $(W/U, V \to W)$
has an object $O_\varphi = (\varphi : V \to U, \text{id} : V \to V)$ for every
$\varphi : V \to U$, and moreover for any object there is a unique
morphism from one of the $O_\varphi$ into it. The result follows.
\end{proof}

\begin{lemma}
\label{lemma-describe-j-shriek-representable}
Let $\mathcal{C}$ be a site.
Let $U \in \Ob(\mathcal{C})$.
Let $X/U$ be an object of $\mathcal{C}/U$.
Then we have $j_{U!}(h_{X/U}^\#) = h_X^\#$.
\end{lemma}

\begin{proof}
Denote $p : X \to U$ the structure morphism of $X$.
By Lemma \ref{lemma-describe-j-shriek} we see $j_{U!}(h_{X/U}^\#)$
is the sheaf associated to the presheaf
$$
V
\longmapsto
\coprod\nolimits_{\varphi \in \Mor_\mathcal{C}(V, U)}
\{\psi : V \to X \mid p \circ \psi = \varphi\}
$$
This is clearly the same thing as $\Mor_\mathcal{C}(V, X)$.
Hence the lemma follows.
\end{proof}

\noindent
We have $j_{U!}(*) = h_U^\#$ by either of the
two lemmas above. Hence for every sheaf
$\mathcal{G}$ over $\mathcal{C}/U$ there is a canonical map
of sheaves $j_{U!}\mathcal{G} \to h_U^\#$. This characterizes
sheaves in the essential image of $j_{U!}$.

\begin{lemma}
\label{lemma-essential-image-j-shriek}
Let $\mathcal{C}$ be a site.
Let $U \in \Ob(\mathcal{C})$.
The functor $j_{U!}$ gives an equivalence of categories
$$
\Sh(\mathcal{C}/U)
\longrightarrow
\Sh(\mathcal{C})/h_U^\#
$$
\end{lemma}

\begin{proof}
Let us denote objects of $\mathcal{C}/U$ as pairs $(X, a)$
where $X$ is an object of $\mathcal{C}$ and $a : X \to U$ is
a morphism of $\mathcal{C}$. Similarly, objects of
$\Sh(\mathcal{C})/h_U^\#$ are pairs $(\mathcal{F}, \varphi)$.
The functor $\Sh(\mathcal{C}/U) \to \Sh(\mathcal{C})/h_U^\#$
sends $\mathcal{G}$ to the pair $(j_{U!}\mathcal{G}, \gamma)$
where $\gamma$ is the composition of
$j_{U!}\mathcal{G} \to j_{U!}*$ with the identification
$j_{U!}* = h_U^\#$.

\medskip\noindent
Let us construct a functor from
$\Sh(\mathcal{C})/h_U^\#$ to $\Sh(\mathcal{C}/U)$.
Suppose that $(\mathcal{F}, \varphi)$ is given.
For an object $(X, a)$ of $\mathcal{C}/U$
we consider the set $\mathcal{F}_\varphi(X, a)$
of elements $s \in \mathcal{F}(X)$ which under $\varphi$ map to the image
of $a \in \Mor_\mathcal{C}(X, U) = h_U(X)$ in
$h_U^\#(X)$. It is easy to see that
$(X, a) \mapsto \mathcal{F}_\varphi(X, a)$ is
a sheaf on $\mathcal{C}/U$. Clearly, the rule
$(\mathcal{F}, \varphi) \mapsto \mathcal{F}_\varphi$
defines a functor $\Sh(\mathcal{C})/h_U^\# \to \Sh(\mathcal{C}/U)$.

\medskip\noindent
Consider also the functor
$\textit{PSh}(\mathcal{C})/h_U \to \textit{PSh}(\mathcal{C}/U)$,
$(\mathcal{F}, \varphi) \mapsto \mathcal{F}_\varphi$
where $\mathcal{F}_\varphi(X, a)$ is defined as the set of elements
of $\mathcal{F}(X)$ mapping to $a \in h_U(X)$.
We claim that the diagram
$$
\xymatrix{
\textit{PSh}(\mathcal{C})/h_U \ar[r] \ar[d] &
\textit{PSh}(\mathcal{C}/U) \ar[d] \\
\Sh(\mathcal{C})/h_U^\# \ar[r] &
\Sh(\mathcal{C}/U)
}
$$
commutes, where the vertical arrows are given by sheafification.
To see this\footnote{An alternative is to describe
$\mathcal{F}_\varphi$ by the cartesian diagram
$$
\vcenter{
\xymatrix{
\mathcal{F}_\varphi \ar[r] \ar[d] & {*} \ar[d] \\
\mathcal{F}|_{\mathcal{C}/U} \ar[r] & h_U|_{\mathcal{C}/U}
}
}
\quad\text{for presheaves and}\quad
\vcenter{
\xymatrix{
\mathcal{F}_\varphi \ar[r] \ar[d] & {*} \ar[d] \\
\mathcal{F}|_{\mathcal{C}/U} \ar[r] & h_U^\#|_{\mathcal{C}/U}
}
}
$$
for sheaves and use that restriction to $\mathcal{C}/U$ commutes
with sheafification.}, it
suffices to prove that the construction commutes with
the functor $\mathcal{F} \mapsto \mathcal{F}^+$ of
Lemmas \ref{lemma-plus-presheaf} and \ref{lemma-plus-functorial}
and Theorem \ref{theorem-plus}.
Commutation with $\mathcal{F} \mapsto \mathcal{F}^+$ follows from the fact
that given $(X, a)$ the categories of coverings of $(X, a)$ in
$\mathcal{C}/U$ and coverings of $X$ in $\mathcal{C}$
are canonically identified.

\medskip\noindent
Next, let $\textit{PSh}(\mathcal{C}/U) \to \textit{PSh}(\mathcal{C})/h_U$
send $\mathcal{G}$ to the pair $(j_{U!}^{PSh}\mathcal{G}, \gamma)$
where $j_{U!}^{PSh}\mathcal{G}$ the presheaf defined by the formula
in Lemma \ref{lemma-describe-j-shriek} and $\gamma$ is the composition of
$j_{U!}^{PSh}\mathcal{G} \to j_{U!}*$ with the identification
$j_{U!}^{PSh}* = h_U$ (obvious from the formula).
Then it is immediately clear that the diagram
$$
\xymatrix{
\textit{PSh}(\mathcal{C}/U) \ar[r] \ar[d] &
\textit{PSh}(\mathcal{C})/h_U \ar[d] \\
\Sh(\mathcal{C}/U) \ar[r] &
\Sh(\mathcal{C})/h_U^\#
}
$$
commutes, where the vertical arrows are sheafification.
Putting everything together it suffices to show there are
functorial isomorphisms $(j_{U!}^{PSh}\mathcal{G})_\gamma = \mathcal{G}$
for $\mathcal{G}$ in $\textit{PSh}(\mathcal{C}/U)$
and $j_{U!}^{PSh}\mathcal{F}_\varphi = \mathcal{F}$
for $(\mathcal{F}, \varphi)$ in $\textit{PSh}(\mathcal{C})/h_U$.
The value of the presheaf $(j_{U!}^{PSh}\mathcal{G})_\gamma$
on $(X, a)$ is the fibre of the map
$$
\coprod\nolimits_{a' : X \to U} \mathcal{G}(X, a') \to \Mor_\mathcal{C}(X, U)
$$
over $a$ which is $\mathcal{G}(X, a)$. This proves the first equality.
The value of the presheaf $j_{U!}^{PSh}\mathcal{F}_\varphi$ is
on $X$ is
$$
\coprod\nolimits_{a : X \to U} \mathcal{F}_\varphi(X, a) =
\mathcal{F}(X)
$$
because given a set map $S \to S'$ the set $S$ is the disjoint
union of its fibres.
\end{proof}

\noindent
Lemma \ref{lemma-essential-image-j-shriek}
says the functor $j_{U!}$ is the composition
$$
\Sh(\mathcal{C}/U) \rightarrow
\Sh(\mathcal{C})/h_U^\# \rightarrow
\Sh(\mathcal{C})
$$
where the first arrow is an equivalence.

\begin{lemma}
\label{lemma-j-shriek-commutes-equalizers-fibre-products}
Let $\mathcal{C}$ be a site. Let $U \in \Ob(\mathcal{C})$.
The functor $j_{U!}$ commutes with fibre products and equalizers (and
more generally finite connected limits). In particular, if
$\mathcal{F} \subset \mathcal{F}'$ in $\Sh(\mathcal{C}/U)$, then
$j_{U!}\mathcal{F} \subset j_{U!}\mathcal{F}'$.
\end{lemma}

\begin{proof}
Via Lemma \ref{lemma-essential-image-j-shriek}
and the fact that an equivalence of categories commutes
with all limits, this reduces to the fact that the functor
$\Sh(\mathcal{C})/h_U^\# \rightarrow \Sh(\mathcal{C})$
commutes with fibre products and equalizers. Alternatively, one can
prove this directly using the description of $j_{U!}$ in
Lemma \ref{lemma-describe-j-shriek}
using that sheafification is exact. (Also, in case $\mathcal{C}$ has
fibre products and equalizers, the result follows from
Lemma \ref{lemma-preserve-equalizers}.)
\end{proof}

\begin{lemma}
\label{lemma-j-shriek-reflects-injections-surjections}
Let $\mathcal{C}$ be a site. Let $U \in \Ob(\mathcal{C})$.
The functor $j_{U!}$ reflects injections and surjections.
\end{lemma}

\begin{proof}
We have to show $j_{U!}$ reflects monomorphisms and epimorphisms, see
Lemma \ref{lemma-mono-epi-sheaves}.
Via Lemma \ref{lemma-essential-image-j-shriek}
this reduces to the fact that the functor
$\Sh(\mathcal{C})/h_U^\# \to \Sh(\mathcal{C})$
reflects monomorphisms and epimorphisms.
\end{proof}

\begin{lemma}
\label{lemma-compute-j-shriek-restrict}
Let $\mathcal{C}$ be a site. Let $U \in \Ob(\mathcal{C})$.
For any sheaf $\mathcal{F}$ on $\mathcal{C}$ we have
$j_{U!}j_U^{-1}\mathcal{F} = \mathcal{F} \times h_U^\#$.
\end{lemma}

\begin{proof}
This is clear from the description of $j_{U!}$ in
Lemma \ref{lemma-describe-j-shriek}.
\end{proof}

\begin{lemma}
\label{lemma-relocalize}
Let $\mathcal{C}$ be a site.
Let $f : V \to U$ be a morphism of $\mathcal{C}$.
Then there exists a commutative diagram
$$
\xymatrix{
\mathcal{C}/V \ar[rd]_{j_V} \ar[rr]_j & &
\mathcal{C}/U \ar[ld]^{j_U} \\
& \mathcal{C} &
}
$$
of continuous and cocontinuous functors.
The functor $j : \mathcal{C}/V \to \mathcal{C}/U$,
$(a : W \to V) \mapsto (f \circ a : W \to U)$
is identified with the functor
$j_{V/U} : (\mathcal{C}/U)/(V/U) \to \mathcal{C}/U$
via the identification $(\mathcal{C}/U)/(V/U) = \mathcal{C}/V$.
Moreover we have $j_{V!} = j_{U!} \circ j_!$,
$j_V^{-1} = j^{-1} \circ j_U^{-1}$, and
$j_{V*} = j_{U*} \circ j_*$.
\end{lemma}

\begin{proof}
The commutativity of the diagram is immediate.
The agreement of $j$ with $j_{V/U}$ follows from the definitions. By
Lemma \ref{lemma-composition-cocontinuous}
we see that the following diagram of morphisms of topoi
\begin{equation}
\label{equation-relocalize}
\vcenter{
\xymatrix{
\Sh(\mathcal{C}/V) \ar[rd]_{j_V} \ar[rr]_j & &
\Sh(\mathcal{C}/U) \ar[ld]^{j_U} \\
& \Sh(\mathcal{C}) &
}
}
\end{equation}
is commutative. This proves that
$j_V^{-1} = j^{-1} \circ j_U^{-1}$ and $j_{V*} = j_{U*} \circ j_*$.
The equality $j_{V!} = j_{U!} \circ j_!$
follows formally from adjointness properties.
\end{proof}

\begin{lemma}
\label{lemma-relocalize-explicit}
Notation $\mathcal{C}$, $f : V \to U$, $j_U$, $j_V$, and $j$ as in
Lemma \ref{lemma-relocalize}. Via the identifications
$\Sh(\mathcal{C}/V) = \Sh(\mathcal{C})/h_V^\#$
and
$\Sh(\mathcal{C}/U) = \Sh(\mathcal{C})/h_U^\#$
of
Lemma \ref{lemma-essential-image-j-shriek}
we have
\begin{enumerate}
\item the functor $j^{-1}$ has the following description
$$
j^{-1}(\mathcal{H} \xrightarrow{\varphi} h_U^\#)
=
(\mathcal{H} \times_{\varphi, h_U^\#, f} h_V^\# \to h_V^\#).
$$
\item the functor $j_!$ has the following description
$$
j_!(\mathcal{H} \xrightarrow{\varphi} h_V^\#) =
(\mathcal{H} \xrightarrow{h_f \circ \varphi} h_U^\#)
$$
\end{enumerate}
\end{lemma}

\begin{proof}
Proof of (2). Recall that the identification
$\Sh(\mathcal{C}/V) \to \Sh(\mathcal{C})/h_V^\#$
sends $\mathcal{G}$ to $j_{V!}\mathcal{G} \to j_{V!}(*) = h_V^\#$
and similarly for
$\Sh(\mathcal{C}/U) \to \Sh(\mathcal{C})/h_U^\#$.
Thus $j_!\mathcal{G}$ is mapped to
$j_{U!}(j_!\mathcal{G}) \to j_{U!}(*) = h_U^\#$
and (2) follows because $j_{U!}j_! = j_{V!}$
by Lemma \ref{lemma-relocalize}.

\medskip\noindent
The reader can now prove (1) by using that $j^{-1}$ is the
right adjoint to $j_!$ and using that the rule in (1)
is the right adjoint to the rule in (2). Here is a direct proof.
Suppose that $\varphi : \mathcal{H} \to h_U^\#$ is an object of
$\Sh(\mathcal{C})/h_U^\#$. By the proof of
Lemma \ref{lemma-essential-image-j-shriek}
this corresponds to the sheaf
$\mathcal{H}_\varphi$ on $\mathcal{C}/U$ defined by the rule
$$
(a : W \to U)
\longmapsto
\{ s \in \mathcal{H}(W) \mid \varphi(s) = a\}
$$
on $\mathcal{C}/U$. The pullback $j^{-1}\mathcal{H}_\varphi$ to
$\mathcal{C}/V$ is given by the rule
$$
(a : W \to V)
\longmapsto
\{ s \in \mathcal{H}(W) \mid \varphi(s) = f \circ a\}
$$
by the description of $j^{-1} = j_{U/V}^{-1}$ as the restriction
of $\mathcal{H}_\varphi$ to $\mathcal{C}/V$.
On the other hand, applying the rule to the object
$$
\xymatrix{
\mathcal{H}' = \mathcal{H} \times_{\varphi, h_U^\#, f} h_V^\#
\ar[rr]^-{\varphi'} & & h_V^\#
}
$$
of $\Sh(\mathcal{C})/h_V^\#$
we get $\mathcal{H}'_{\varphi'}$
given by
\begin{align*}
(a : W \to V)
\longmapsto
&  \{ s' \in \mathcal{H}'(W) \mid \varphi'(s') = a\} \\
= &
\{ (s, a') \in \mathcal{H}(W) \times h_V^\#(W) \mid
a' = a \text{ and } \varphi(s) = f \circ a'\}
\end{align*}
which is exactly the same rule as the one describing
$j^{-1}\mathcal{H}_\varphi$ above.
\end{proof}

\begin{remark}
\label{remark-localize-presheaves}
Localization and presheaves. Let $\mathcal{C}$ be a category.
Let $U$ be an object of $\mathcal{C}$. Strictly speaking the functors
$j_U^{-1}$, $j_{U*}$ and $j_{U!}$ have not been defined for presheaves.
But of course, we can think of a presheaf as a sheaf for the
chaotic topology on $\mathcal{C}$ (see Example \ref{example-indiscrete}).
Hence we also obtain a functor
$$
j_U^{-1} :
\textit{PSh}(\mathcal{C})
\longrightarrow
\textit{PSh}(\mathcal{C}/U)
$$
and functors
$$
j_{U*}, j_{U!} :
\textit{PSh}(\mathcal{C}/U)
\longrightarrow
\textit{PSh}(\mathcal{C})
$$
which are right, left adjoint to $j_U^{-1}$. By
Lemma \ref{lemma-describe-j-shriek}
we see that $j_{U!}\mathcal{G}$ is the presheaf
$$
V \longmapsto
\coprod\nolimits_{\varphi \in \Mor_\mathcal{C}(V, U)}
\mathcal{G}(V \xrightarrow{\varphi} U)
$$
In addition the functor $j_{U!}$ commutes with fibre products and
equalizers.
\end{remark}

\begin{remark}
\label{remark-localization-cartesian-cocontinuous}
Let $\mathcal{C}$ be a site. Let $U \to V$ be a morphism of $\mathcal{C}$.
The cocontinuous functors $\mathcal{C}/U \to \mathcal{C}$ and
$j : \mathcal{C}/U \to \mathcal{C}/V$ (Lemma \ref{lemma-relocalize})
satisfy property $P$ of Remark \ref{remark-cartesian-cocontinuous}.
For example, if we have objects $(X/U)$, $(W/V)$, a morphism
$g : j(X/U) \to (W/V)$, and a covering $\{f_i  : (W_i/V) \to (W/V)\}$ then
$(X \times_W W_i/U)$ is an avatar of $(X/U) \times_{g, (W/V), f_i} (W_i/V)$
and the family $\{(X \times_W W_i/U) \to (X/U)\}$ is a covering
of $\mathcal{C}/U$.
\end{remark}






